\chapter{Modules, imports and macros}
\label{ch:modules}

\section{Modules}
\label{sec:modules}
\index{modules}\index{mod@\texttt{mod}}\index{use@\texttt{use}}

\fc|mod Name where ... end| groups definitions, types, traits and nested
modules under a name. Outside the module, its items are written with a
qualified name, \fc|Name.item|; inside, they are used directly.
\fc|use Name| brings a module's items into scope for the rest of the file.
Modules nest: \fc|A.B.x|, \fc|use A.B|.

\forgeexample{modules}

\section{Imports}
\label{sec:imports}
\index{import@\texttt{import}}
\fc|import "path"| inserts the definitions of another file at that point.
The path is relative to the importing file, and the \texttt{.forge}
extension may be omitted. A file imported twice is included once, and an
import cycle is an error.

\forgeexample*{import-lib}
\forgeexample{import-main}

\section{Macros}
\label{sec:macros}
\index{macros}\index{macro_rules@\texttt{macro\_rules"!}}

\fc|macro_rules!| defines a macro by example, in the style of Rust. A macro
has one or more arms, separated by semicolons; each arm is a \emph{matcher}
and a \emph{transcriber}. An invocation \fc|name!(...)|, \fc|name![...]| or
\fc|name!{...}| is replaced, before type checking, by the transcriber of the
first arm whose matcher fits its arguments.

\begin{itemize}
\item In a matcher, \fc|$x:kind| captures a fragment of the given kind:
  \texttt{expr}, \texttt{pat}, \texttt{ty} (or \texttt{type}),
  \texttt{ident}, \texttt{literal} (or \texttt{lit}), \texttt{path},
  \texttt{block}, \texttt{item} or \texttt{tt} (a single token tree). Other
  tokens must appear literally.
\item \fc|$( ... )*| repeats a part zero or more times, \fc|$( ... )+| one
  or more times; a separator can come before the \fc|*| or \fc|+|:
  \fc|$($x:expr),*|.
\item An \texttt{expr} fragment is substituted as a single expression, as if
  parenthesised: with \fc|$e * $e|, \fc|square!(1 + 2)| is \fc|9|.
\end{itemize}

The built-in macros are \fc|count_args!(...)| (the number of
comma-separated arguments), \fc|vec!(...)| (a list), \fc|concat!(...)| (the
concatenation of string literals) and \fc|stringify!(...)| (a textual dump of
the argument's syntax tree).

\forgeexample{macros}

Macros are expanded before type checking, so a macro can expand to any
expression; they cannot be defined inside a block.
